\stagechapter{D}{Material capture}

\noindent\textbf{What this stage covers.} What is collected for each app, in what form, and how it is made citable later: store pages and structured store disclosures, privacy policies and terms, provider pages, screens and recordings from use, and the metadata (URL, date, language, version, hash) that lets a finding be traced to its source. The draft thesis plan has a core corpus per app (policy, store descriptions, structured store disclosures, the main product page and privacy or security pages linked from it or the policy), a fixed supplementary search for provider blog posts, opportunistic ads, German as primary language with English fallback, and an archive record (URL, date, language, file hash, app version). Open decisions: D2 (the collector and its settings), D5, D7, D11 (device scope: apps only).

\section{Approach 1: the environment of expected use and the capture of context}
\pa{light2018} splits the walkthrough method into the \emph{environment of expected use} and the technical walkthrough. The first part records the app's vision (purpose, target users, scenarios of use) from store descriptions, websites, blogs and press material \pc{light2018}{9}{An app’s vision involves its purpose, target user base and scenarios of use}, the operating model (price, in-app purchase menus, who is paid in data) \pc{light2018}{10}{The app’s price in app stores, in-app purchase menus} and governance (terms of service, FAQs, hard-coded rules), studied over time with the Internet Archive \pc{light2018}{11}{conceptions of governance, especially when studied over time using tools like the Internet Archive}. This is the source for treating the store page, website and policy as one corpus, and for recording versions with the Wayback Machine. \pa{duguay2023} adds that a clean account is not neutral: record the persona and account configuration, app version, operating system and date, and note interface changes during collection, since jurisdiction-dependent interfaces matter for a Germany-specific sample.

\section{Approach 2: structured store disclosures as a separate surface}
Store labels are self-reported presentations. \pa{magoulas2026} recorded the full Data Safety schema for 49{,}578 apps (data shared, data collected, security practices; 14 data categories) and distinguishes \emph{not completed} from \emph{declared none} and \emph{declared negative}; \pa{li2022} captured every US App Store privacy label weekly for 32 weeks with an error-checking and re-crawl step and fixed denominators; \pa{arkalakis2024} captured Data Safety sections longitudinally and compared them with app behaviour, finding a large gap between label and behaviour; \pa{malki2024} extracted the Data Safety fields manually for 20 apps and searched the same fields in the policy by keyword. For a communication-only thesis the label is a claim to be recorded with its capture date, never ground truth. \pa{kelley2009} is the conceptual ancestor of these labels and gives criteria (fixed rows, explicit ``not collected'', no ambiguous symbols) for judging them.

\section{Approach 3: policy capture and versioning}
\begin{itemize}
\item \pa{shipp2020} hashed captured policy versions; \pa{lin2025} archived policies as text with time stamps; \pa{doan2025} used dated policy capture for six companies; \pa{sunyaev2015} fixed the discovery protocol (store, developer site, first 30 web results).
\item Scope matters at capture: one policy may cover several apps or only a website (\pa{malki2024}: the shortest policy covered only website cookies; one developer shared a generic policy across apps). The draft plan codes shared policies once and links them to all relevant apps.
\item Terms of service and caveats sit next to policies (\pa{robillard2019} includes terms of agreement). \pa{wp292018} says the privacy information of an app must be specific to that app, available before download and within two taps, so the capture should record \emph{where} each document was reached from.
\item \pa{bfarm2023}: a change in the location of data storage is a notifiable change for a DiGA, which is one more reason to keep dated policy versions and to expect updates during the study.
\end{itemize}

\section{Approach 4: material from use}
\pa{huckvale2015} froze the software state during evaluation (no app or OS updates), used simulated plausible data with test flags and ran an exhaustive sweep of user-interface elements before realistic use. \pa{malki2024} wiped and reinstalled every app to simulate first use; \pa{alfawzan2022} used an emulator (Android 7.1.2) and an iPhone 11 (iOS 14.4) and recorded the device; \pa{lin2025} fixed a time budget per app from a pilot (22.4 minutes on average, 15--30 minutes allowed). \pa{light2018} asks for detailed field notes, screenshots, screen video and recorded thoughts \pc{light2018}{11}{producing a collection of data by generating detailed field notes and recordings}, and stresses that the order of screens is data too \pc{light2018}{13}{Paying attention not only to the app’s features but also to the flow of activity}. \pa{habib2019} uses a fixed template for each capture step, study-only accounts and a clean browser for each site.

\section{Approach 5: scale, automation and its limits}
Automated collection brings reach and loses depth: \pa{tangari2021} (static analysis, traffic capture and review mining on 20{,}991 apps) and \pa{li2022} are census designs; \pa{shipp2020} and \pa{huckvale2019} add traffic capture to policy reading. The thesis excludes network analysis, so what it records are \emph{statements and interfaces}, not behaviour. \pa{aoir2019} flags that platforms may prohibit automated collection, that the legal status for researchers is unsettled, and that data minimisation, storage and retention apply to scraped store data; the thesis should state which store terms apply to its collectors and how much is stored.

\begin{table}[t]
\centering\small
\caption{Material captured per surface and the metadata that makes it reusable.}
\label{tab:D-material}
\begin{tabularx}{\textwidth}{@{}P{2.7cm}YY@{}}
\toprule
\textbf{Surface} & \textbf{Minimum record} & \textbf{Precedent}\\
\midrule
Store listing & text, screenshots, category, version, last update, ratings and counts, price, in-app purchases, date, storefront & \pa{magoulas2026}, \pa{li2022}, \pa{apa2024}, \pa{stoyanov2015mars}\\
Structured disclosure & every field of the label, with ``not completed'' distinguished from ``none'' & \pa{magoulas2026}, \pa{arkalakis2024}\\
Policy and terms & full text, URL, date, language, hash, how reached (store, app, website), scope (one or several apps) & \pa{shipp2020}, \pa{sunyaev2015}, \pa{lin2025}, \pa{wp292018}\\
Provider pages & privacy and security pages linked from product page or policy; blog search with fixed window & \pa{light2018}, \pa{huckvale2015}\\
Use material & screenshots or video per step, device, OS, app version, account state, synthetic data, time stamps & \pa{light2018}, \pa{huckvale2015}, \pa{malki2024}, \pa{habib2019}\\
Directory (DiGA) & listing status, manufacturer, data-processing locations, confirmed checklists & \pa{bfarm2023}\\
\bottomrule
\end{tabularx}
\end{table}

\begin{protocolbox}[Capture checklist assembled from the corpus]
\begin{itemize}[leftmargin=1.3em]
\item Record storefront, country, language, device, OS, account state and persona (\pa{duguay2023}, \pa{rethlefsen2021}).
\item Save raw responses unchanged, then derived tables; hash each document; store capture date and time (\pa{shipp2020}).
\item Log how each document was reached and whether the link was broken or empty (\pa{tangari2021}, \pa{sunyaev2015}).
\item Freeze app versions during the walkthrough window and log updates that arrive anyway (\pa{huckvale2015}).
\item Keep the original wording; add translations separately and note the language per source (D7).
\item Keep personal data out of the corpus; check store terms before automated retrieval (\pa{aoir2019}).
\end{itemize}
\end{protocolbox}

\begin{thesisbox}[Implications for this thesis (stage D)]
The draft plan already follows the best precedents (archive record with hash and version, language per source, shared policies coded once). Points to settle: (1) treat Data Safety and App Store labels as claims with capture dates and separate denominators for ``not completed''; (2) decide how many policy versions to keep during the study window and how to detect changes (\pa{li2022}'s snapshot cadence is a model); (3) record the persona and account configuration for walkthroughs, and the apps' regional behaviour; (4) write a short collector note stating store terms, rate and retention (D2, \pa{aoir2019}); (5) report absence only as ``no mention identified in the captured material'', as the draft already specifies.
\end{thesisbox}
